\section{\ENG{LMS--SCORM Engine Integration Architecture}}
\label{sec:lms_scorm_integration}

Η διασύνδεση ανάμεσα στο \ENG{LMS} και στο υπόλοιπο οικοσύστημα υλοποιείται σε δύο διακριτές ροές. Η πρώτη είναι περισσότερο διαχειριστική και αφορά την εισαγωγή μαθημάτων, τη δημιουργία ή τον συγχρονισμό χρηστών και τη διαχείριση εγγραφών. Σε αυτή τη ροή, το \ENG{LMS} καλεί το \ENG{engine} μέσω του \ENG{SDK}, χρησιμοποιώντας \ENG{engine-level} \ENG{JWT} με ρόλους \ENG{ADMIN} ή \ENG{INSTRUCTOR}. Η δεύτερη είναι η \ENG{run-time} ροή του εκπαιδευόμενου. Εκεί το \ENG{LMS} ζητά τη δημιουργία \ENG{launch} και λαμβάνει \texttt{\ENG{launchUrl}} και \texttt{\ENG{launchToken}}. Από το σημείο αυτό και μετά, η πλοήγηση μεταφέρεται στον \ENG{player}, ο οποίος λειτουργεί ως \ENG{browser-facing boundary service}.

Το κρίσιμο αρχιτεκτονικό χαρακτηριστικό είναι ότι ο \ENG{player} ζητά από το \ENG{engine} το \ENG{launch context} μέσω του \texttt{\ENG{GET /api/v1/launches/\{launchId\}}}, λαμβάνει πληροφορία για τον χρήστη, το μάθημα, το \texttt{\ENG{standard}}, το \texttt{\ENG{entrypoint}} και την πηγή του περιεχομένου, και μόνο τότε συνθέτει τη σελίδα εκκίνησης. Η πληροφορία του περιεχομένου επιστρέφεται ως \ENG{presigned MinIO ZIP URL}, γεγονός που αποδεικνύει ότι η διανομή του πακέτου και η εκτέλεσή του είναι ξεχωριστή φάση από την ίδια τη δημιουργία του \ENG{launch}.

Το \ENG{LMS} παραμένει υπεύθυνο για τοπικά στοιχεία \ENG{BL (Business Logic)}, όπως η πλοήγηση της διεπαφής, οι συσχετίσεις τοπικών χρηστών με \ENG{engine users}, και η τελική προβολή της προόδου. Δεν αναλαμβάνει όμως ούτε τη διαχείριση \ENG{SCORM API calls} ούτε την αποθήκευση της πρώτης ύλης του \ENG{runtime}. Αντιστρόφως, ο \ENG{engine} δεν ενδιαφέρεται για το πώς παρουσιάζεται η διοικητική εμπειρία του \ENG{LMS}, αλλά για το αν το σύνολο \ENG{course-user-enrollment-attempt-launch} είναι συνεπές και ασφαλώς εξουσιοδοτημένο.


Το Σχήμα~\ref{fig:lms_scorm_integration} αποτυπώνει τη διπλή αυτή ροή ολοκλήρωσης. Ο Πίνακας~\ref{tab:repository_role_mapping} συμπυκνώνει τη χαρτογράφηση των βασικών αποθετηρίων στα αντίστοιχα αρχιτεκτονικά καθήκοντα, ώστε τα επόμενα κεφάλαια να μπορούν να αναφέρονται σε συγκεκριμένα \ENG{implementation boundaries} χωρίς να επαναλαμβάνεται ο συνολικός χάρτης της πλατφόρμας.

\begin{figure}[H]
  \centering
\begin{chapterfigurebox}
  \begin{tikzpicture}[
      font=\footnotesize,
      node distance=1.9cm and 2.5cm,
      box/.style={draw, rounded corners, fill=gray!5, align=center, text width=3.3cm, minimum height=1.1cm},
      emph/.style={draw, rounded corners, fill=gray!12, align=center, text width=3.6cm, minimum height=1.1cm, font=\footnotesize\bfseries},
      store/.style={draw, cylinder, shape border rotate=90, aspect=0.28, minimum height=1.2cm, minimum width=3.2cm, align=center, fill=gray!5},
      arrow/.style={-Latex, thick},
      boundary/.style={draw, dashed, rounded corners, inner sep=11pt}
    ]
    \node[box] (admin) {\ENG{Admin / Instructor}\\\ENG{in LMS}};
    \node[box, below=of admin] (lms) {\ENG{Laravel LMS}};
    \node[box, right=of lms] (sdk) {\ENG{PHP SDK}};
    \node[emph, right=of sdk] (engine) {\ENG{SCORM Engine}};

    \node[box, below=2.0cm of lms] (browser) {\ENG{Learner}\\\ENG{Browser}};
    \node[emph, below=2.0cm of engine] (player) {\ENG{SCORM Player}};
    \node[store, below=1.9cm of player] (minio) {\ENG{MinIO}\\\ENG{course package}};

    \draw[arrow] (admin) -- node[left]{\ENG{admin actions}} (lms);
    \draw[arrow] (lms) -- node[above, align=center]{\ENG{typed API}\\\ENG{calls}} (sdk);
    \draw[arrow] (sdk) -- node[above, align=center]{\ENG{import / users}\\\ENG{/ enrollments / launches}} (engine);
    \draw[arrow] (lms) -- node[left, align=center]{\ENG{launch}\\\ENG{URL}} (browser);
    \draw[arrow] (browser) -- node[below]{\ENG{open} \texttt{\ENG{/launch}}} (player);
    \draw[arrow] (player.north) to[bend left=14] node[right, align=center]{\ENG{launch}\\\ENG{context}} (engine.south);
    \draw[arrow] (player) -- node[right, align=center]{\ENG{download}\\\ENG{ZIP}} (minio);
    \draw[arrow] (player.north west) to[bend right=18] node[left, align=center]{\ENG{commit}\\\ENG{/ terminate}} (engine.south west);

    \node[boundary, fit=(admin)(lms)(browser)] (clientboundary) {};
    \node[font=\footnotesize, above=0.12cm of clientboundary] {\ENG{LMS-owned experience}};

    \node[boundary, fit=(engine)(player)(minio)] (runtimeboundary) {};
    \node[font=\footnotesize, above=0.12cm of runtimeboundary] {\ENG{Runtime and content delivery perimeter}};
  \end{tikzpicture}
\end{chapterfigurebox}
  \caption{Αρχιτεκτονική ολοκλήρωσης της πλατφόρμας με διάκριση διοικητικής και \ENG{learner-facing} ροής.}
  \label{fig:lms_scorm_integration}
\end{figure}


\begin{table}[H]
\centering
\caption{Χαρτογράφηση αποθετηρίων και αρχιτεκτονικών ρόλων του υλοποιημένου οικοσυστήματος}
\label{tab:repository_role_mapping}
\begingroup
\footnotesize
\setlength{\tabcolsep}{3pt}
\renewcommand{\arraystretch}{1.20}
\begin{adjustbox}{center,max width=\textwidth}
\begin{tabular}{|p{3.0cm}|p{4.1cm}|p{6.5cm}|}
\hline
\ENG{Repository} & Κύριος ρόλος & Ισχυρότερα \ENG{implementation} τεκμήρια \\
\hline
\texttt{\ENG{example-lms-client}} & \ENG{LMS} διεπαφή, διοικητικές ροές, \ENG{learner pages} & \texttt{\ENG{ScormEngineGateway}}, \texttt{\ENG{LearnerController}}, \texttt{\ENG{launch.blade.php}}, \texttt{\ENG{progress.blade.php}} \\
\hline
\texttt{\ENG{scorm-engine-php-sdk}} & Τυποποιημένο \ENG{integration layer} προς το \ENG{engine} & \texttt{\ENG{CoursesApi}}, \texttt{\ENG{LaunchesApi}}, \texttt{\ENG{AttemptsApi}}, \texttt{\ENG{DtoMapper}} \\
\hline
\texttt{\ENG{scorm-engine}} & Κεντρικός \ENG{backend} πυρήνας για \ENG{import}, \ENG{launch}, εκτέλεση, μόνιμη αποθήκευση και \ENG{reporting} & \texttt{\ENG{CourseController}}, \texttt{\ENG{LaunchController}}, \texttt{\ENG{RuntimeController}}, \ENG{orchestrators}, \ENG{migrations} \\
\hline
\texttt{\ENG{player}} & \ENG{Browser-facing} υπηρεσία εκκίνησης, διανομής περιεχομένου και προώθησης \ENG{commit/terminate} & \texttt{\ENG{launchRoute}}, \texttt{\ENG{internalRuntimeRoute}}, \texttt{\ENG{RenderLaunchCommandHandler}}, \texttt{\ENG{LaunchPageRenderer}} \\
\hline
\texttt{\ENG{central-docker-infrastructure}} & \ENG{Deployment composition}, \ENG{stores}, \ENG{observability stack} και λειτουργικά \ENG{checks} & \texttt{\ENG{docker-compose.yml}}, \texttt{\ENG{smoke-check.sh}}, \texttt{\ENG{validate-env.sh}}, \texttt{\ENG{logstash.conf}} \\
\hline
\end{tabular}
\end{adjustbox}
\endgroup

\end{table}
